% Source: https://artofproblemsolving.com/community/c3h1989118_2019_amc_10a_problem_13
% Stable identity: aops-topic-1989118-q01
Since $BC$ is a diameter, the angles inscribed in the semicircle are right angles: $\angle BDC=90^\circ$ and $\angle BEC=90^\circ$. So $BD\perp AC$ and $CE\perp AB$; the diagonals $BD$ and $CE$ of $BCDE$ are two altitudes of $\triangle ABC$, and $F$ is its orthocenter.

The base angles are $\angle CAB=\angle CBA=\dfrac{180^\circ-40^\circ}{2}=70^\circ$. In quadrilateral $AEFD$ the angles at $E$ and $D$ are right angles, so
\[
\angle EFD=360^\circ-90^\circ-90^\circ-70^\circ=110^\circ,
\]
and $\angle BFC=\angle EFD$ (vertical angles).

\begin{asy}
usepackage("amsmath");
size(5cm);
defaultpen(fontsize(10pt));
real s=Sin(20), h=Cos(20);
pair C=(0,h), A=(-s,0), B=(s,0);
pair M=(B+C)/2; real r=abs(B-C)/2;
pair footp(pair P, pair U, pair V){ return U+dot(P-U,V-U)/dot(V-U,V-U)*(V-U); }
pair D=footp(B,A,C), Ee=footp(C,A,B);
pair F=extension(B,D,C,Ee);
draw(circle(M,r),gray(0.55));
draw(A--B--C--cycle,linewidth(1));
draw(B--D,blue); draw(C--Ee,blue); draw(D--Ee,dashed+gray(0.4));
dot(A);dot(B);dot(C);dot(D);dot(Ee);dot(F);
label("$A$",A,SW);label("$B$",B,SE);label("$C$",C,N);label("$D$",D,W);label("$E$",Ee,S);label("$F$",F,dir(200));
draw(arc(C,0.13,250,290));label("$40^\circ$",C+0.13*dir(250),W,fontsize(8pt));
label("\mbox{}",(min(currentpicture,user=true)+max(currentpicture,user=true))/2,black); // keeps the text after the figure black when LaTeX typesets the labels inline
\end{asy}

So $\angle BFC=\boxed{110^\circ}$, choice $\textbf{(D)}$.
